\chapterpage{profiles}{从热图提取截面}
\subhead{沿鼠标画出的线段取值}
按第~\ref{heatmap}~章绘好示例热图。在热图页打开\ui{工具 → 热图分布提取}（\keys{Ctrl+H}），工具窗口会保持打开，主画布仍可操作。
\begin{enumerate}
\item 点击\ui{开始提取}。在热图约 50 µm 高度，从最早时刻向最晚时刻按住左键水平拖动，松手后会弹出分布曲线窗口。
\item 在结果窗口选择“导出数据”或“保存当前图片”。关闭结果窗口后，可以在提取工具中点击“停止”结束取样，点击“清除标记”去掉热图上的线段。
\end{enumerate}
\Figure{profile}{沿 50 µm 附近水平取样，曲线峰值出现在这条线段的中部。手工拖动时，端点稍有偏差就会改变峰值读数。}{86mm}
程序沿线段从头至尾等距取 200 个点，并在已显示的热图网格上插值。结果横轴 \ui{Distance along line} 表示从起点算起的坐标距离。本例沿水平线取样，只有时间发生变化，距离数值等于经过的 ns 数。若斜着取线，距离同时使用两轴的数值差。保存曲线时，应记下两个端点及其单位，便于他人还原取样位置。
\checknote{精确指定位置}{要在 50 µm 的准确位置、按原始输出时刻取值，可使用下一页的“固定位置”截面。它与鼠标线段提取使用不同的采样点，导出表的列也不同。}

\nextpage{extrema}{在指定截面上寻找极值}
\begin{enumerate}
\item 保持热图页为当前页面，打开\ui{工具 → 极值寻找工具}（\keys{Ctrl+E}）。选“固定位置”，输入 \code{50}，点击“添加Profile”。用英文逗号分隔可一次添加多个位置。
\item 保留全部范围，勾选“寻找极大值”，点击\ui{提取并寻找极值}。切换到“分布曲线”。本例在 2 ns 找到 8 eV 的最大值。
\end{enumerate}
\Figure{extrema}{红色标记给出固定 50 µm 截面上的最大值。左侧表格列出了本次选定的位置和搜索区间。}{86mm}
\subhead{固定位置与固定时间}
固定位置时，程序在每个已存储时刻沿空间插值到指定位置。固定时间时，程序在文件覆盖的时间范围内选择最接近输入值的一条时间记录，并使用该步自身的坐标。若需要精确复现某一帧，直接填写实际输出时刻。
在分布表的范围栏中可以限定搜索区间。例如对 50 µm 的时间曲线输入 \code{1-3}，即只在 1 至 3 ns 内找极值。每个区间可分别求最大值或最小值。“导数曲线”使用数值梯度。峰较宽、靠近边界或中间有缺测时，应同时查看原始采样点。可以保存图像或导出截面数值。
\checknote{导出时区分单位}{极值工具的 CSV 首列为 \code{X}，其余每列对应一个截面。建议将固定位置和固定时间的截面分别导出，使每个文件的首列具有一致的物理单位。}

\chapterpage{annotations}{标注图形与检查诊断结果}
\subhead{给报告中的曲线加注释}
显示三条空间曲线后按 \keys{Ctrl+M}。在标注工具中选择“箭头”，设定颜色和线宽，从准备放文字的位置拖向 2 ns 的峰。改选“文字”，点击箭头旁边，在弹窗中输入 \code{2 ns}。要删除某一个标记，可以在标注列表中选中该项再删除。
\Figure{annotation}{箭头和文字标出了中间的温度峰。保存带标注图像，可以让读者直接识别该曲线及其对应时刻。}{72mm}
工具还支持直线、矩形和椭圆，并可设置透明度、线型及文字样式。请在标注仍可见时使用“导出带标注图片”。重新绘制会重建坐标轴，建议先保存标注图像，再修改绘图选项。标签页导出的数值表记录原始取值。
\subhead{查看检测工具给出的轨迹}
在适用的数据集上，先切换到热图，再打开\ui{工具 → 冲击波检测}（\keys{Ctrl+D}）。窗口可显示场量、剖面和检测轨迹，并可保存热图、导出轨迹数据。逐个时刻核对轨迹点与原始剖面上的位置，同时保留所用场量、坐标、单位和配置，方便复查。练习包只包含平滑移动的温度峰，适合演示前面的绘图与取样操作。
\checknote{恢复画布操作}{关闭标注工具或停止线段提取后，再用左键拖动进行普通框选缩放。}

\chapterpage{export}{保存图像与导出数值}
在当前绘图页按 \keys{Ctrl+S} 保存图像。PNG 适合直接插入报告，PDF 适合排版。输入文件名时写上扩展名。主绘图页保存栅格图像的默认分辨率是 300 dpi。
\Table{命令所在位置}{导出文件的内容}{
空间图 → 导出数据 & 第一列是坐标，随后每列对应一个已选且启用的时刻，表头中的时间使用 ns。\\
时间图 → 导出数据 & 第一列是所选单位的时间。全局模式随后是一列标量，网格模式随后每列对应一个已选且启用的单元。\\
热图 → 导出数据 & 普通空间场量按 500 个位置重新取样，首列为坐标，其余每列为一个已记录时刻。多群边界谱另按频率群导出。\\
线段结果窗口 → 导出数据 & 两列分别为 \code{Distance} 和变量名，共 200 个采样点。请另记轴方向、线段端点和单位。\\
极值寻找工具 → 导出数据 & 首列为共享的 \code{X}，其余列对应各截面。某截面缺少对应横坐标时，单元格可能留空。\\}
\subhead{顶栏“导出”菜单的两种输出}
“导出为 CSV”最多取可用列表中的前五个空间变量，每个变量单独写一份表。列表示时刻，行按数组顺序排列，数值采用解析器中的存储单位。表中没有另加位置列。输入例如 \code{run.csv} 的小写扩展名，生成的文件名会包含变量，如 \code{run_T.csv}。
“导出为 ASCII”调用旧式写出器。输入 \code{run.txt}，将分别生成 \code{run_SCout.txt} 和 \code{run_TSout.txt}。旧式写出器需要预设的一组变量。练习数据等精简文件若缺少这些字段，请在相应绘图页导出 CSV。
\checknote{交付数据时}{在可写目录中为每次导出使用清楚的文件名，同时记录源文件、变量、坐标、单位、所选时刻或单元，以及用过的插值和坐标范围。}

\nextpage{checkcsv}{核验导出的表格}
保留一份原始 \code{quickstart.zip}。另建输出目录，在空间图、网格时间图和热图中分别保存 \code{profiles.csv}、\code{cell.csv} 和 \code{map.csv}。按 UTF-8、逗号分隔格式导入电子表格。下表的行数均仅计算数据行。
\Table{文件}{练习数据的检查项}{
\code{profiles.csv} & 选 \#0、\#4、\#8 后有 21 行、4 列。在 50 µm 这一行，2 ns 对应列应为 8.000000e+00 eV。\\
\code{cell.csv} & 单独选 \#10 后有 9 行、2 列。第 5 行数据的时间为 2.000000e+00 ns，同一行的温度为 8.000000e+00 eV。\\
\code{map.csv} & 有 500 行位置、10 列数据（含坐标列）。第一列从 0 到 100 µm。500 为偶数，中间两行分别位于 50 µm 两侧。\\
全局时间曲线 CSV & 9 行、2 列。表头为 \code{Time (ns)} 和 \code{ENIE}。2 ns 时标量的存储数值为 14。\\}
\subhead{空间图 CSV 的实际表头}
\begin{verbatim}
XC Centroid (um),t=0.0000ns T (eV),t=2.0000ns T (eV),
t=4.0000ns T (eV)
\end{verbatim}
这里为便于阅读做了折行，文件中四个字段在同一行。导入时把第一行设为列名，复制或另存数据时保留科学计数法。
\subhead{运动网格导出时的坐标}
在 2.3.1 版中，空间图 CSV 的唯一坐标列取自第一条时间记录。普通热图的 CSV 也使用这个参考网格对每一时刻的场量重新取样。屏幕上的欧拉热图则使用各时刻自己的位置。对运动网格做定量比较时，应从源输出中取同一时刻的场量和 \code{XC} 或 \code{X} 配对。若要比较某个固定空间位置在不同时刻的值，可用极值工具的“固定位置”截面，并保留坐标单位。
\checknote{为什么会差一点}{固定时刻的空间曲线保留原有的单元采样。热图网格点和手画线上的采样点都可能落在两个中心之间。核对之前，先统一时间、位置和数值单位。}

\chapterpage{recorder}{录制并回放绘图操作}
练习使用 \code{quickstart.out}。生成的 Python 脚本内嵌标准二进制解析器。GUI 还可以读取十进制算例和支持的格式化输入，这些文件的图像和数值表可直接从 GUI 中导出。
\begin{enumerate}
\item 打开 \code{quickstart.out}。按 \keys{Ctrl+R} 打开“记录代码”，点击\ui{开始录制}。当时已打开的文件也会计入操作数。
\item 回到空间图，选择 \ui{T}、\ui{XC}、\ui{um}、\ui{eV} 和时刻 \ui{\#4}，点击\ui{绘制}。回到录制窗口，点击\ui{停止录制}，再点\ui{生成脚本}。
\item 将脚本命名为 \code{example_replay.py}，保存到自己可写的目录，并把 \code{quickstart.out} 放在旁边。在装有 Python、NumPy、SciPy、Matplotlib 的终端中切换到该目录，运行下列命令。
\end{enumerate}
\begin{verbatim}
python example_replay.py
\end{verbatim}
\Figure{recorder}{先开文件再录一张图，停止后计数为 2。两项记录分别对应已载入的数据和一次绘图操作，此时可以生成脚本。}{54mm}
示例会在当前工作目录生成 \code{example_replay_fig1_spatial_plot.png}。脚本先尝试录制时的绝对输入路径，找不到时再检查脚本旁边的同名文件。打开结果图，应看到 2 ns 的一条曲线，峰位在 50 µm 附近。若要核查准确的单元数值，另存空间图 CSV。
\checknote{开始下一次录制}{点击“开始录制”会清空先前的操作列表。先为需要保留的操作命名并生成脚本。录制器记录支持的绘图操作，交互式标注请单独保存成图片。}

\chapterpage{trouble}{常见问题与诊断}
\Table{遇到的问题}{检查方法}{
双击后看到包内文件 & 在“打开方式”中选择软件安装器，或执行第~\ref{install}~章的本地 APT／DNF 安装命令。查看包管理器给出的依赖错误。\\
菜单启动后没有窗口 & 在图形桌面的终端运行 \code{multi1dvisual}，查看第一条 Qt 或缺包信息。远程终端需要能够连接图形显示。\\
中文字符显示为空框 & 用发行版的字体管理工具安装中文字体后重启程序。Debian／Ubuntu 安装包推荐 \code{fonts-noto-cjk}。\\
十进制文件打不开 & 先解压练习包。核对两个必需表的同名文件前缀及大小写、制表符、\code{TIME} 表头以及各时间块的行数和顺序。\\
曲线缺失或看似水平 & 检查当前文件、变量、选择列表和排除状态。恢复坐标范围自动，关闭对数轴后重新绘制，再看原始列名和单位。\\
热图留白或整片同色 & 选择线性色标并点击颜色“自动”，检查轴范围，在空间图查看原始记录是否有有限数值。超出原始位置覆盖范围的区域可能留白。\\
工具提示需要热图 & 切至热图，选择有数据的场量并实际绘制，再打开工具。拖动行为异常时，先结束旧的标注或提取操作。\\
帮助菜单打不开手册 & 到安装目录确认 PDF 是否存在，尝试用本机 PDF 阅读器直接打开，再查看桌面默认的 PDF 打开方式。\\}
\subhead{查找最近一次日志}
打开用户主目录，显示隐藏文件夹，进入 \textsf{.MULTI1D可视化工具} 下的 \code{logs}。每次运行会写入以 \code{run_} 开头的日志，程序保留最近三份。要让他人复现问题，请附上软件版本、操作系统、输入格式、所选页面与单位、触发问题的步骤，以及最新日志中的相关几行。

\chapterpage{reference}{快捷键与文件速查}
\Table{快捷键}{作用}{
\keys{Ctrl+O}／\keys{Ctrl+S}／\keys{Ctrl+Q} & 打开输入文件／保存当前图像／退出。\\
\keys{Ctrl+1}／\keys{Ctrl+2}／\keys{Ctrl+3} & 依次切换到空间图、时间图或热图。\\
\keys{Ctrl+M} & 对当前画布打开标注工具。\\
\keys{Ctrl+H}／\keys{Ctrl+E} & 在热图中打开线段分布提取／极值寻找工具。\\
\keys{Ctrl+D}／\keys{Ctrl+R} & 从热图打开冲击波检测／打开代码录制窗口。\\
\keys{A}／\keys{H} & 焦点位于画布时，恢复完整视图／回退一级缩放。\\}
\subhead{常用变量名}
\Table{文件中的键}{含义}{
\code{TIME}、\code{XC}、\code{X} & 分别为输出时间、中心位置和界面位置。显示换算前，时间使用 s，位置使用 cm。\\
\code{T}、\code{TI}、\code{TR} & 分别为电子温度、离子温度和辐射温度。练习数据提供 \code{T}。\\
\code{R}、\code{P}、\code{PT} & 分别为质量密度、压强和总压强。坐标轴标注取决于变量和所选显示单位。\\
\code{ENIE}、\code{ENII}、\code{ENKI} & 分别为全局电子内能、离子内能和动能。若图中未标注能量单位，以输出程序对该字段的定义为准。\\}
\subhead{偏好设置与项目文件}
上一页的用户目录中还保存 \code{config.json} 和 \code{language_config.json}。项目首页列出两种语言的手册和练习数据。练习压缩包内含 \code{quickstart.case}、\code{quickstart.center.dat}、\code{quickstart.scalars.dat} 和 \code{quickstart.out}。如果要查看合成数值的具体定义，源码中有 \code{docs/examples/make_quickstart.py}。
\checknote{项目}{软件作者为徐博。仓库中提供源码、安装包构建说明和两份手册的可编辑源文件。\repo}
